% Generated draft: requires mathematical and source review.
\documentclass[11pt]{article}
\usepackage[margin=25mm]{geometry}
\usepackage{fontspec}
\setmainfont{Latin Modern Roman}
\usepackage{amsmath,amssymb,mathtools,unicode-math}
\setmathfont{Latin Modern Math}
\usepackage{enumitem,needspace,xurl,hyperref}
\setlength{\parindent}{0pt}
\setlength{\parskip}{5pt}
\setlength{\emergencystretch}{3em}
\newcommand{\N}{\mathbb{N}}
\newcommand{\Z}{\mathbb{Z}}
\newcommand{\Q}{\mathbb{Q}}
\newcommand{\R}{\mathbb{R}}
\newcommand{\C}{\mathbb{C}}
\newcommand{\F}{\mathbb{F}}
\newcommand{\A}{\mathbb{A}}
\newcommand{\PP}{\mathbb{P}}
\newcommand{\E}{\mathbb{E}}
\newcommand{\eps}{\varepsilon}
\newcommand{\dd}{\,\mathrm{d}}
\begin{document}
\section*{UnsolvedMath ID: 30001220}
\textbf{Source problem code:} catalogue entry
\subsection{Definitions and mathematical statement}
The supplied record concerns a Hecke algebra of type \(B_n\) with unequal parameters \(Q,q\) and a weight parameter \(\xi\in\R_{>0}\setminus\Q\). The rank range, coefficient ring, quadratic relations, and relation of \(\xi\) to \(Q,q\) are not specified. Consequently, the following describes the requested assertion as a formulation schema; it does not resolve those missing choices.

For a positive integer \(n\), a bipartition of \(n\) is an ordered pair \(\lambda=(\lambda^{(1)},\lambda^{(2)})\) of partitions satisfying \(|\lambda^{(1)}|+|\lambda^{(2)}|=n\). Write \(\Lambda_n\) for their set. A standard bitableau of shape \(\lambda\) fills the two Young diagrams bijectively with \(1,\ldots,n\), with entries increasing along each row and down each column within each diagram. Denote their set by \(\operatorname{Std}(\lambda)\).

After fixing a commutative coefficient ring \(A\) with identity and the intended Hecke algebra \(H=H_A(B_n;Q,q)\), the requested cellular datum consists of a partial order on \(\Lambda_n\), an \(A\)-basis
\[
\{C^\lambda_{s,t}:\lambda\in\Lambda_n,\ s,t\in\operatorname{Std}(\lambda)\},
\]
and an \(A\)-linear algebra anti-involution \(*\) satisfying \((C^\lambda_{s,t})^*=C^\lambda_{t,s}\). With \(H^{>\lambda}\) denoting the \(A\)-span of basis elements whose labels are strictly greater than \(\lambda\), the cellular multiplication condition is that, for every \(h\in H\),
\[
hC^\lambda_{s,t}\equiv
\sum_{u\in\operatorname{Std}(\lambda)}r_h(u,s)C^\lambda_{u,t}
\pmod{H^{>\lambda}},
\]
where \(r_h(u,s)\in A\) is independent of \(t\). The resulting cell module \(S_A^\lambda\) has basis indexed by \(\operatorname{Std}(\lambda)\), with action given by these coefficients. The source asks that these be the intended Specht modules and that the datum arise from Kazhdan--Lusztig theory; it does not specify the proposed basis, order, anti-involution, or identification with Specht modules.

For a permitted specialization \(\theta:A\to K\) to a field, put \(H_K=K\otimes_{A,\theta}H\) and \(S_K^\lambda=K\otimes_{A,\theta}S_A^\lambda\). Using the specialized cellular bilinear form, put
\[
D_K^\lambda=S_K^\lambda/\operatorname{rad}S_K^\lambda.
\]
The requested assertion is that the Kazhdan--Lusztig construction supplies the above datum, that its nonzero simple quotients realize all of Ariki's parametrizations, and that the decomposition numbers
\[
d_{\lambda\mu}=[S_K^\lambda:D_K^\mu],\qquad D_K^\mu\ne0,
\]
agree with the intended canonical-basis transition matrix under the intended correspondence of labels.

The source does not identify the permitted specializations, the family meant by ``all of Ariki's parametrizations,'' or the vector space, two bases, evaluation, and indexing needed to specify that transition matrix. It also does not say whether realizing all parametrizations requires varying \(\xi\). Thus a fully quantified matrix identity and a complete unconditional conjecture cannot be recovered from this record alone. In particular, it remains unclear whether Lusztig's conjectures P1 through P15 are hypotheses of the requested construction or part of what must be established.
\subsection{Short English statement}
The question asks whether unequal-parameter Kazhdan--Lusztig theory for type \(B\), with a positive irrational weight, produces a cellular basis indexed by bipartitions and standard bitableaux, giving the intended Specht modules, Ariki parametrizations, and canonical-basis decomposition numbers. Essential parameter, specialization, and indexing conventions are missing, so this is a qualified statement-only draft requiring mathematical review.
\subsection{Sources}
\begin{itemize}
\item Local UnsolvedMath registry: supplied record with permanent ID \(30001220\), source problem code \texttt{catalogue entry}, title ``Cellular Models for Unequal Parameter Hecke Algebras,'' in the Oberwolfach Reports Open Problems collection. The supplied proposed year is 2009; no individual proposer is identified.
\item Supplied report citation: \emph{Representations of Finite Groups} (2009), Oberwolfach Reports, pp.\ 917--970, DOI \texttt{10.4171/owr/2009/17}. Source URL and literature-source label ``original report'': \url{https://doi.org/10.4171/owr/2009/17}. The supplied background describes these reports as documenting mathematical workshops and their open-problem discussions.
\item Additional supplied literature source, labeled ``rank three results'': \url{https://arxiv.org/abs/1906.09451}. The record supplies no authors, title, or precise result for this reference.
\item Literature assessment dated exactly \texttt{2026-08-21}: the supplied status is ``partially solved,'' with the assessment that the unequal-parameter construction remains conditional in general although substantial cases of Lusztig's P1 through P15 are proved. This is supplied dated curation derived from the OWR material, not an independently verified literature assessment. The record explicitly requests independent verification of time-sensitive status and cited claims.
\item The registry labels the statement \texttt{corrected\_verified}, explaining that the complete statement was preserved while publication formatting was normalized. This supplied label does not establish mathematical correctness or independent source verification of this draft.
\end{itemize}
\end{document}
